\section{Projektmanagement und Projektreflexion}

\subsection{Projektmanagement}

\subsubsection{Projektphasen}

Das Projekt wurde in klar abgegrenzte Phasen unterteilt, um
Analyse, Konzeption, Umsetzung und Validierung strukturiert
durchführen zu können. Die Übergänge zwischen den Phasen waren
teilweise iterativ, da Erkenntnisse aus der Implementierung
wieder in die Ausgestaltung des SOLL-Prozesses eingeflossen
sind.

\begin{table}[H]
\centering
\begin{tabular}{p{5cm}p{4cm}p{3cm}}
\toprule
Phase & Zeitraum & Status \\
\midrule
Analyse und Anforderungsabstimmung & xx.xx - xx.xx & abgeschlossen \\
Architektur und Prozessmodellierung & xx.xx - xx.xx & abgeschlossen \\
Implementierung & xx.xx - xx.xx & abgeschlossen \\
Integration und Test & xx.xx - xx.xx & abgeschlossen \\
Abschluss und Ergebnisaufbereitung & xx.xx - xx.xx & abgeschlossen \\
\bottomrule
\end{tabular}
\end{table}

\subsubsection{Projektplan}

Der Projektplan diente als gemeinsame Orientierung für
inhaltliche Meilensteine, Abstimmungstermine und die
Reihenfolge der Umsetzung. Besonders hilfreich war die frühe
Sichtbarkeit von Abhängigkeiten zwischen Prozessmodell,
Schnittstellen und Benutzeroberfläche.

\begin{figure}[H]
\centering
\includegraphics[width=\textwidth]{../Projektplan/assets/projectplan.png}
\caption{Projektplan}
\end{figure}

\subsubsection{Projektaufwand}

Der dokumentierte Aufwand verteilt sich auf Analyse,
Konzeption, Implementierung, Test und Abstimmung mit dem
Auftraggeber. Die Tabelle dient als Beispielstruktur zur
Erfassung der geleisteten Stunden und kann mit den tatsächlichen
Projektdaten vervollständigt werden.

\begin{longtable}{p{2cm}p{3cm}p{5cm}p{2cm}}
\toprule
Datum & Person & Tätigkeit & Stunden \\
\midrule
xx.xx.2026 & Name 1 & Anforderungsanalyse und Recherche & 4 \\
\midrule
xx.xx.2026 & Name 2 & Backend-Implementierung & 6 \\
\midrule
xx.xx.2026 & Name 3 & BPMN-Modellierung und Camunda-Integration & 5 \\
\midrule
xx.xx.2026 & Name 1 & Frontend-Umsetzung und UI-Anpassungen & 5 \\
\midrule
xx.xx.2026 & Team & Integrationstest und Ergebnisabstimmung & 4 \\
\bottomrule
\end{longtable}

\subsection{Projektreflexion}

\subsubsection{Herausforderungen}

Herausfordernd war insbesondere die Abstimmung an den
Schnittstellen zwischen fachlicher Prozessbeschreibung und
technischer Umsetzung. Bereits kleine Unklarheiten bei
Statusübergängen oder Rückmeldewegen hatten direkte
Auswirkungen auf mehrere Systemkomponenten.

Auch die Konsistenz zwischen Benutzerführung, Datenmodell und
Workflowlogik erforderte wiederholte Abstimmungen.

\subsubsection{Lessons Learned}

Für vergleichbare Projekte empfiehlt es sich, Statusmodelle,
Ereignisse und Verantwortlichkeiten frühzeitig gemeinsam zu
definieren und schriftlich festzuhalten. Darüber hinaus hat
sich gezeigt, dass ein visuelles Prozessmodell in Kombination
mit einer frühen technischen Teilimplementierung sehr hilfreich
ist, um Missverständnisse schnell sichtbar zu machen.
